\appendix
\section{Explicit estimates and certificate templates}\label{sec:appendices}

This appendix records details that are routine in classical notation but
important for an iPRA implementation: every ``eventually'', ``uniformly'', and
``interchange the order'' is replaced by a bounded arithmetic inequality.

\subsection{A uniform hypergeometric ratio bound}

Let
\[
 T_k=\frac{\prod_{i=1}^{r}(\alpha_i)_k}
           {\prod_{j=1}^{s}(\beta_j)_k}\frac{z^k}{k!}.
\]
Suppose rational bounds $\abs{\alpha_i}\le B$ and
$\abs{\beta_j}\le B$ are certified and every $\beta_j+k$ is apart from zero.
For $k\ge2B+1$,
\[
 \abs{k+\alpha_i}\le\frac32k,
 \qquad
 \abs{k+\beta_j}\ge\frac12k,
 \qquad k+1\ge k.
\]
Consequently
\begin{equation}\label{eq:appendix-ratio}
 \abs{\frac{T_{k+1}}{T_k}}
 \le \abs z\,\frac{3^r2^{s-r}}{k^{s+1-r}}.
\end{equation}
For $r\le s$, bounded search finds $N$ making the right side at most $1/2$;
then
\begin{equation}\label{eq:appendix-tail-half}
 \abs{\sum_{k=N}^{\infty}T_k}\le2\abs{T_N}.
\end{equation}
For $r=s+1$ and $\abs z\le1-\epsilon$, use the sharper expansion
\[
 \frac{\prod_i(k+\alpha_i)}{(k+1)\prod_j(k+\beta_j)}
 =1+O(1/k)
\]
with a rational constant $C$.  Choosing
$N\ge2C/\epsilon$ gives ratio at most $1-\epsilon/2$ and tail
$2\abs{T_N}/\epsilon$.  The $O(1/k)$ assertion is implemented by expanding the
two monic integer-cleared polynomials and bounding their lower coefficients.
No asymptotic notation appears in the checker.

For $r>s+1$, the nonterminating power series generally has zero radius.  The
report never calls this series convergent; the corresponding $G$-function is a
contour/residue summation name.  If an upper parameter is a nonpositive integer,
the series terminates and is a finite polynomial, covered by
\cref{thm:polynomial-repair}.

\subsection{Derivation of the right Slater coefficient}

Put $s=b_h+k+u$.  Since
\[
 \Gamma(-k-u)=\frac{(-1)^{k+1}}{k!u}+O(1),
\]
the ordinary residue of the integrand has a leading minus sign.  The right
loop is negatively oriented, so the contour integral contributes the positive
coefficient $(-1)^k/k!$.  The remaining factors are rewritten using
\begin{align*}
 \Gamma(c+k)&=\Gamma(c)(c)_k,\\
 \Gamma(c-k)&=(-1)^k\frac{\Gamma(c)}{(1-c)_k},\\
 \frac1{\Gamma(c-k)}&=(-1)^k\frac{(1-c)_k}{\Gamma(c)}.
\end{align*}
After collecting signs, the argument is
$(-1)^{p-m-n}z$ and the constant term is
\begin{equation}\label{eq:appendix-Ch}
 C_h=
 \frac{\displaystyle
   \prod_{\substack{1\le j\le m\\j\ne h}}\Gamma(b_j-b_h)
   \prod_{j=1}^{n}\Gamma(1+b_h-a_j)}
 {\displaystyle
   \prod_{j=m+1}^{q}\Gamma(1+b_h-b_j)
   \prod_{j=n+1}^{p}\Gamma(a_j-b_h)}.
\end{equation}
The $k$th coefficient is therefore
\[
 C_hz^{b_h}
 \frac{\prod_{j=1}^{p}(1+b_h-a_j)_k}
      {\prod_{\substack{1\le j\le q\\j\ne h}}
       (1+b_h-b_j)_k}
 \frac{[(-1)^{p-m-n}z]^k}{k!},
\]
which is exactly the $h$th hypergeometric term in
\cref{thm:slater-right}.  Every equation above has a denominator-cleared
factorial proof for rational parameters away from poles; for general certified
complex parameters it is a finite Gamma recurrence ball equality.

The left formula follows by setting $s=a_h-1-k+u$.  Here the left loop has
positive orientation, and the same three Gamma identities produce the dual
${}_qF_{p-1}$ expression.

\subsection{Confluent residues without parameter limits}

Suppose $r$ Gamma factors have a pole at $s_0$.  Write
\[
 F(s_0+u)=\sum_{j=-r}^{r-1}c_ju^j+u^rR(u),
 \qquad \abs{R(u)}\le M
 \quad(\abs u\le\rho).
\]
Multiplying by
\[
 z^{s_0+u}=z^{s_0}\sum_{j=0}^{r-1}
 \frac{(u\Log z)^j}{j!}+u^rR_z(u)
\]
shows that the residue is
\begin{equation}\label{eq:appendix-confluent-res}
 z^{s_0}\sum_{j=0}^{r-1}c_{-1-j}\frac{(\Log z)^j}{j!}.
\end{equation}
Thus the degree in $\Log z$ is at most $r-1$.  To certify the remainder, use
truncated convolution.  If
$A=A_{<r}+u^rR_A$ and $B=B_{<r}+u^rR_B$ on $\abs u\le\rho$, then
\[
 \abs{R_{AB}}
 \le \abs{A_{<r}}_\rho M_B+
      \abs{B_{<r}}_\rho M_A+\rho^rM_AM_B,
\]
where $\abs{\sum a_ju^j}_\rho=\sum\abs{a_j}\rho^j$ after shifting away
negative powers.  This formula is a primitive-recursive update of rational
bounds.  It is the formal replacement for phrases such as ``take the limit as
$b_2\to b_1$''.

\subsection{The differential equation as a polynomial identity}

Let $\thetaop=z\frac{\dd}{\dd z}$.  Acting on $z^s$ replaces $\thetaop$ by
$s$.  Therefore the differential operator in \cref{thm:ode} multiplies
the kernel by
\[
 P(s)=(-1)^{p-m-n}z\prod_{j=1}^{p}(s-a_j+1)
      -\prod_{j=1}^{q}(s-b_j).
\]
Gamma recurrence gives the exact shift identity
\begin{equation}\label{eq:appendix-kernel-shift}
 \Phi(s-1)=(-1)^{p-m-n}
 \frac{\prod_{j=1}^{p}(s-a_j)}
      {\prod_{j=1}^{q}(s-b_j-1)}\Phi(s),
\end{equation}
with the convention of \eqref{eq:main-definition}.  After replacing $s$ by
$s+1$ and multiplying by $z^{s+1}$, the first term in $P(s)$ is the shifted
second term.  On a finite contour the substitution is exact; the two short
connector segments cancel or tend below the allocated error by the same
vertical estimate.  In a residue series, \eqref{eq:appendix-kernel-shift} is
precisely the coefficient recurrence.  Hence the ODE is proved without a
classical theorem about differentiating arbitrary improper integrals.

\subsection{Why the nonresonance correction is necessary}

For $p\le q$, the $h$th proposed local solution begins with
$C_hz^{b_h}(1+O(z))$.  The first $q$ Euler derivatives at a positive rational
base point have leading matrix
\[
 V=(b_h^{\ell})_{0\le\ell<q,\ 1\le h\le q}.
\]
Its determinant is the Vandermonde product
$\prod_{i<j}(b_j-b_i)$.  A quantitative lower bound for this product proves
linear independence and supplies a conditioning constant.  If two exponents
coincide, or differ by an integer in a regular-singular equation, this argument
fails and logarithmic solutions may occur.  The Laurent-jet construction
produces those solutions.  Therefore the unqualified word ``fundamental'' on
the source page cannot be accepted literally in resonant cases.

\subsection{A finite Fubini budget}

Consider a formally desired exchange
\[
 \int_0^\infty\int_{-\infty}^{\infty}F(x,t)\dd t\dd x
 =\int_{-\infty}^{\infty}\int_0^\infty F(x,t)\dd x\dd t.
\]
At precision $2^{-p}$ choose rational $0<\epsilon<R$ and $T>0$ so that
\begin{align*}
 \int_0^\epsilon\int_{\R}\abs F&\le2^{-p-4},&
 \int_R^\infty\int_{\R}\abs F&\le2^{-p-4},\\
 \int_\epsilon^R\int_{\abs t>T}\abs F&\le2^{-p-4},&
 E_{\rm quad}&\le2^{-p-4}.
\end{align*}
On the finite rectangle $[\epsilon,R]\times[-T,T]$, both iterated integrals
are represented by the same finite double quadrature sum, so reordering is
literal associativity/commutativity.  The four discarded errors total at most
$2^{-p-2}$; the remaining budget covers the inner Gamma/Beta evaluation and
rounding.  This is the pattern used in every integral theorem of
\cref{sec:transforms}.

\subsection{Branch transport by overlapping boxes}

A branch certificate is not merely a symbolic choice of $\Arg$.  It contains a
rational box $D$, a rational center $z_0$, a value $L_0$ approximating
$\Log z_0$, and a margin $d>0$ separating $D$ from zero and the cut.  If
$z_1$ lies in the same box and
$\abs{(z_1-z_0)/(z_1+z_0)}\le r<1$, then
\begin{equation}\label{eq:appendix-log-step}
 \Log z_1-\Log z_0
 =2\sum_{k=0}^{N-1}\frac{w^{2k+1}}{2k+1}+R_N,
 \quad w=\frac{z_1-z_0}{z_1+z_0},
\end{equation}
with
\[
 \abs{R_N}\le\frac{2r^{2N+1}}{(2N+1)(1-r^2)}.
\]
A finite chain of overlapping boxes transports the branch.  Powers are then
$z^a=\exp(a\Log z)$ with a propagated ball.  This construction proves every
branch continuation used in the report and makes monodromy visible: a loop
whose certified logarithm increment is $2\pi\ii$ changes $z^a$ by
$\e^{2\pi\ii a}$ rather than pretending the function is globally
single-valued.

\subsection{Representative iPRA sentence schemas}

For a fixed compiled cosine specification, the mathematical statement becomes
\[
 \forall x\,\forall p\,
 \bigl(\Safe_{\cos}(x,p)\to
 \BallEq(\sqrt\pi\,G(x^2/4),\cos(x),p)\bigr).
\]
The displayed quantifiers are ordinary first-order quantifiers over natural
codes.  The functions \code{SafeCos}, \code{meijerCos}, \code{cosEval}, and
\code{ballEq} are fixed primitive-recursive symbols.  The actual Lean target is
a closed formula built from their graph equations and a proof
\code{[] |- [PRAAx] phi}.

For Mellin transformation, a fixed certificate producer $C$ gives
\[
 \forall u\,\forall p\,
 \bigl(\operatorname{MellinCert}(C,u,p)\to
 \BallEq(\operatorname{MellinEval}(u,p),
         \operatorname{GammaRatio}(u,p),p)\bigr).
\]
This is not second-order quantification over functions or certificates.  The
programs are compiled into the language before the theorem is stated.

\subsection{Reproducibility and limitations}

The archive contains:
\begin{enumerate}
\item the full LaTeX source and compiled PDF;
\item a CSV ledger and its generated typeset table;
\item a Lean-oriented host scaffold and explicit object-language targets;
\item exact coefficient checks using rational arithmetic;
\item high-precision numerical regression checks using \code{mpmath};
\item source provenance, build logs, rendered PDF page images, and checksums.
\end{enumerate}

The executable checks do not prove the universal analytic theorems.  Conversely,
the human proofs do not imply that a Lean file has been compiled.  The missing
formal work is precisely the construction of the object-language derivations
described in \cref{sec:formalization}.  This boundary is part of the result,
not a hidden caveat.
